\documentclass[../../../include/open-logic-section]{subfiles}

\begin{document}

\olfileid[pt]{sth}{story}{rus}
\olsection{Paradoxo de Russell (de novo)}

Na \olref[sfr][][]{part}, trabalhamos com uma teoria ingênua dos conjuntos. Mas,
segundo uma concepção \emph{muito} ingênua, os conjuntos são apenas as
extensões dos predicados. Este pensamento ingênuo imporia o
seguinte princípio:

\begin{defish}
  \emph{Compreensão ingênua.} $\Setabs{x}{\phi(x)}$ existe para qualquer fórmula $\phi$.
\end{defish}

Por tentador que seja este princípio, é provadamente inconsistente. Vimo-lo na \olref[sfr][set][rus]{sec}, mas o resultado é tão importante, e tão direto, que vale a pena repetir. Ipsis verbis.

\begin{thm}[Paradoxo de Russell]
Não existe conjunto $R = \Setabs{x}{x \notin x}$
\end{thm}

\begin{proof}
Se $R = \Setabs{x}{x \notin x}$ existisse, então
$R \in R$ sse $R \notin R$, o que é uma contradição.
\end{proof}

Russell descobriu este resultado em junho de 1901. (Não colocou, porém,
o paradoxo exatamente na forma em que o acabamos de apresentar, pois estava
a considerar a teoria dos conjuntos de Frege, tal como delineada em \emph{Grundgesetze}. Voltaremos a isto na \olref[blv]{sec}.) Russell escreveu a
Frege a 16 de junho de 1902, explicando a inconsistência no sistema de
Frege. Para a correspondência, e um pouco de contexto, ver
\citet[pp.~124--8]{Heijenoort1967}. 

Vale a pena enfatizar que esta prova de duas linhas é um resultado de
\emph{lógica pura}. De fato, usamos implicitamente um axioma (não lógico?) de extensionalidade na nossa notação $\Setabs{x}{x \notin x}$; pois $\Setabs{x}{\phi(x)}$ deve ser \emph{o único} (pela extensionalidade) conjunto dos $\phi$s, se algum existir. Mas podemos evitar até o vestígio da extensionalidade, bastando enunciar o resultado assim:
\emph{não existe conjunto cujos membros sejam exatamente os conjuntos que não são membros de si próprios}. E isto pouco tem a ver com conjuntos. Como o próprio Russell observou, raciocínio rigorosamente análogo
leva-nos a concluir: \emph{nenhum homem barbeia exatamente os homens que
não se barbeiam a si próprios}. Ou: \emph{nenhum pug cheira exatamente os pugs que
não se cheiram a si próprios}. E assim por diante. Esquematicamente, a forma do
resultado é apenas: 
\[
\lnot \exists x \forall z(Rzx \liff \lnot R zz).
\]
E isto não passa de um teorema (esquema) da lógica de primeira ordem. Consequentemente,
não podemos evitar o paradoxo de Russell apenas mexendo na nossa teoria
dos conjuntos; ele surge antes sequer de \emph{chegarmos} à teoria dos conjuntos. Se formos
usar lógica (clássica) de primeira ordem, temos simplesmente de
\emph{aceitar} que não existe conjunto $R = \Setabs{x}{x\notin x}$. 

A consequência é esta. Se quisermos aceitar a compreensão ingênua
enquanto \emph{evitamos} a inconsistência, não basta mexer na
\emph{teoria dos conjuntos}. Teríamos, em vez disso, de reformular a nossa
\emph{lógica}.

É claro que já foram apresentadas teorias dos conjuntos com lógicas não clássicas.
Mas são---no mínimo---não standard. A abordagem usual
ao paradoxo de Russell é tratá-lo como uma prova direta de não existência
e depois tentar aprender a viver com ela. É essa a
abordagem que seguiremos.

\end{document}
